% Generated draft: requires mathematical and source review.
\documentclass[11pt]{article}
\usepackage[margin=25mm]{geometry}
\usepackage{fontspec}
\setmainfont{Latin Modern Roman}
\usepackage{amsmath,amssymb,mathtools,unicode-math}
\setmathfont{Latin Modern Math}
\usepackage{enumitem,needspace,xurl,hyperref}
\setlength{\parindent}{0pt}
\setlength{\parskip}{5pt}
\setlength{\emergencystretch}{3em}
\newcommand{\N}{\mathbb{N}}
\newcommand{\Z}{\mathbb{Z}}
\newcommand{\Q}{\mathbb{Q}}
\newcommand{\R}{\mathbb{R}}
\newcommand{\C}{\mathbb{C}}
\newcommand{\F}{\mathbb{F}}
\newcommand{\A}{\mathbb{A}}
\newcommand{\PP}{\mathbb{P}}
\newcommand{\E}{\mathbb{E}}
\newcommand{\eps}{\varepsilon}
\newcommand{\dd}{\,\mathrm{d}}
\begin{document}
\section*{UnsolvedMath ID: 20001392}
\textbf{Source problem code:} catalogue entry
\subsection{Definitions and mathematical statement}
Problem 7 asks whether stable laws exist for nonstationary processes and whether a diffusion rate can be found when there is no limiting law. Nonstationarity alone specifies neither a dependence class nor normalization; ``diffusion rate'' might mean growth of a second moment, a logarithmic spreading exponent, or a scale/quantile when variance is infinite. The question therefore needs such choices before it becomes a universal theorem. The supplied aid gives existence examples and separates distributional shape from scale growth.

Fix \(0<\alpha\le2\), and let \(Z_1,Z_2,\ldots\) be independent standard symmetric \(\alpha\)-stable variables, with characteristic function \(\E e^{itZ_k}=e^{-|t|^\alpha}\). For arbitrary deterministic real weights \(c_k\), put
\[
S_n=\sum_{k=1}^nc_kZ_k,\qquad A_n=\left(\sum_{k=1}^n|c_k|^\alpha\right)^{1/\alpha}.
\]
Whenever \(A_n>0\), the aid gives the exact identity in distribution
\[
S_n/A_n\overset{d}=Z_\alpha.
\]
Varying weights can give non-identically distributed increments, so stable normalized laws are compatible with nonstationarity. No asymptotic approximation is needed in this independent stable family. If all weights through time \(n\) vanish, \(S_n=0\) and normalization by \(A_n=0\) is not defined.

Conversely, fix \(0<a<b<\infty\). The aid gives independent Gaussian increments \(X_k\sim N(0,v_k)\), \(v_k\in\{a,b\}\), arranged in alternating dominating variance blocks so that, with \(S_n=\sum_{k=1}^nX_k\) and \(V_n=\sum_{k=1}^nv_k\), the cluster set of \(V_n/n\) is exactly \([a,b]\). Consequently
\[
\left\{\text{weak subsequential limits of }S_n/\sqrt n\right\}
=\{N(0,q):q\in[a,b]\}.
\]
In particular there is no single limiting law under \(\sqrt n\) normalization. Nevertheless
\[
a n\le V_n\le b n,\qquad
\lim_{n\to\infty}\frac{\log\sqrt{\E S_n^2}}{\log n}=\frac12.
\]
Thus a root-mean-square diffusion exponent exists while the normalized distributional shape does not converge. The coefficient \(V_n/n\) itself does not converge, and normalizing instead by \(\sqrt{V_n}\) gives exactly \(N(0,1)\) at each time. This example is not a claim that no normalization can yield a limit.

The aid also reports a sparse weighted stable example of the opposite kind: \(S_n/A_n\) retains the same exact stable law at every time, but
\[
\liminf_{n\to\infty}\frac{\log A_n}{\log n}
<\limsup_{n\to\infty}\frac{\log A_n}{\log n}.
\]
Its summary does not give the weights or numerical exponents. For \(\alpha<2\), variance is infinite, so \(A_n\) is an intrinsic stable scale, not a root-mean-square width. These paired examples show that a scale exponent and a limiting normalized law are separate properties; they do not give stable-limit theorems for arbitrary dependent nonstationary dynamics.

The supplied verification says the grammatical phrase ``there are is'' occurs in the primary PDF and is not an extraction error; the question here merely repairs its grammar. The registry labels the entry partially solved; attempt 1 is a partial result, problem partially solved at the run, internal rigor proved and novelty confidence low. The research aid is machine generated, not peer reviewed, and requires independent verification before citation.
\subsection{Short English statement}
Nonstationarity can coexist with an exact stable normalized law. The supplied examples also separate scale from shape: Gaussian sums can have diffusion exponent \(1/2\) without a limit under \(\sqrt n\) normalization, while weighted stable sums can keep one normalized law even though their logarithmic scale exponent fails to exist. The meaning of diffusion rate and normalization must be specified.
\subsection{Sources}
\begin{itemize}
\item Local UnsolvedMath registry, permanent ID \(20001392\), source code \texttt{catalogue entry}; \texttt{aim-dynamical-systems-notes.json}, \texttt{notes[49]}.
\item AIM, \emph{Stochastic methods for non-equilibrium dynamical systems}, Problem 7: \url{https://aimath.org/pastworkshops/noneqdynsysproblems.pdf}. The supplied verification reports inspection of the three-page PDF and identifies the grammatical defect as present in print; this draft does not claim independent inspection.
\item Supplied research attempt 1 and registry update August 14, 2026. The original literature field is empty and no separate authoritative citation for the examples is supplied.
\end{itemize}
\end{document}
